\documentclass[11pt]{article}
\usepackage[T1]{fontenc}
\usepackage[utf8]{inputenc}
\usepackage{lmodern}
\usepackage[a4paper,margin=1in]{geometry}
\usepackage{microtype}
\usepackage{amsmath,amssymb,amsthm,mathtools}
\usepackage{booktabs,array}
\usepackage[dvipsnames]{xcolor}
\usepackage[colorlinks=true,linkcolor=MidnightBlue,citecolor=MidnightBlue,
            urlcolor=MidnightBlue]{hyperref}
\usepackage{enumitem}
\usepackage{setspace}
\onehalfspacing
\setlist[itemize]{topsep=4pt,itemsep=2pt,leftmargin=1.5em}
\setlist[enumerate]{topsep=4pt,itemsep=2pt,leftmargin=1.8em}

\newtheorem{theorem}{Theorem}[section]
\newtheorem{lemma}[theorem]{Lemma}
\newtheorem{proposition}[theorem]{Proposition}
\newtheorem{corollary}[theorem]{Corollary}
\theoremstyle{definition}
\newtheorem{definition}[theorem]{Definition}
\newtheorem{construction}[theorem]{Construction}
\newtheorem{probe}[theorem]{Probe}
\theoremstyle{remark}
\newtheorem{remark}[theorem]{Remark}
\newtheorem{observation}[theorem]{Observation}

\newcommand{\DR}{\mathsf{DR}}
\newcommand{\PO}{\mathsf{PO}}
\newcommand{\PP}{\mathsf{PP}}
\newcommand{\PPI}{\mathsf{PPI}}
\newcommand{\EQ}{\mathsf{EQ}}
\newcommand{\RCC}{\{\DR,\PO,\PP,\PPI\}}
\newcommand{\RCCall}{\{\DR,\PO,\EQ,\PP,\PPI\}}
\newcommand{\cl}{\operatorname{cl}}
\newcommand{\Tp}{\operatorname{Tp}}
\newcommand{\tp}{\operatorname{tp}}
\newcommand{\comp}{\operatorname{comp}}
\newcommand{\inv}{\operatorname{inv}}
\newcommand{\TypeSafe}{\operatorname{TypeSafe}}
\newcommand{\Need}{\operatorname{Need}}
\newcommand{\Safe}{\operatorname{Safe}}
\newcommand{\Dom}{\mathsf{Dom}}
\newcommand{\ALCIRCC}[1]{\mathcal{ALCI}_{\mathsf{RCC#1}}}

\title{\textbf{Patchwork Preservation under $\ALCIRCC{5}$ Augmentation}\\[0.4em]
\large A probe toward undecidability}
\author{Claude (Anthropic, Opus 4.7)\\
\small Prompted by Michael Wessel}
\date{April 2026}

\begin{document}
\maketitle

\begin{center}
\fbox{\parbox{0.92\textwidth}{%
\textbf{Disclaimer.}\;
This document was generated by Claude (Anthropic), prompted by Michael
Wessel.  The results have not been independently verified.  We invite
scrutiny and corrections.}}
\end{center}

\begin{abstract}
The decidability argument for $\ALCIRCC{5}$ rests on a specific
corollary of the Renz--Nebel patchwork property for atomic
$\mathsf{RCC5}$ networks: that typed arc-consistency is complete for
realizability in the presence of an $\mathcal{ALCI}$ TBox.  If this
\emph{typed patchwork} fails---if there exists a TBox whose
arc-consistent typed networks are not globally realizable---then
undecidability could plausibly hide in the failure gap.  This paper
probes that gap.  We formalize typed patchwork, separate it from
atomic patchwork, attempt three candidate counterexamples, and report
on what we find.  Our finding is that typed patchwork \emph{appears to
hold} in each probe, with every apparent obstruction reducing either
to local arc-inconsistency or to type-level unsatisfiability.  We
identify precisely where the analysis would need to break for
undecidability to emerge, and argue this aligns structurally with the
completeness-extraction paper's ``realized relations are automatically
relation-safe'' key lemma.  The probe is negative, but it sharpens the
question: any future undecidability proof for $\ALCIRCC{5}$ must
either exhibit a typed-patchwork failure of a specific combinatorial
shape, or bypass patchwork entirely.
\end{abstract}

\section{Introduction}
\label{sec:intro}

Every decidability argument for $\ALCIRCC{5}$ in this project---the
quasimodel method, the cover-tree tableau, the split-forest
semantics, the completeness-extraction paper---ultimately appeals to
Renz and Nebel's patchwork property~\cite{RenzNebel1999} for the
$\mathsf{RCC5}$ calculus.  The patchwork property says that a finite
atomic $\mathsf{RCC5}$ network is globally consistent iff it is
path-consistent.  It is the single load-bearing combinatorial fact
behind everything we have claimed.

The question this paper asks is whether the patchwork property
\emph{survives} augmentation with an $\mathcal{ALCI}$ TBox.  More
precisely: does path-consistency of a typed $\mathsf{RCC5}$ network---a
network whose nodes carry Hintikka types and whose edges are locally
type-safe---imply the existence of a global model satisfying the
TBox?  We call this the \emph{typed patchwork} question.

If typed patchwork holds, the decidability argument for
$\ALCIRCC{5}$ is robust to TBox pressure and the currently-claimed
positive result stands.  If typed patchwork \emph{fails}, there is an
opening for undecidability: the failure would mean that some typed
network is locally consistent at every arc but has no global model,
and one could then try to amplify such a failure into an undecidable
reduction.

The reason this is the cleanest question to attack is that both
outcomes are informative.  A failure of typed patchwork does not
automatically give undecidability (one would still need to encode
computation into failure patterns), but it would invalidate every
existing decidability proof in this repository.  A confirmation of
typed patchwork, conversely, explains why all the undecidability
reductions documented in~\cite{README} are blocked and points future
undecidability work toward the one remaining loophole: reductions
that do not pass through any patchwork-style argument at all.

\paragraph{Contributions.}
\begin{itemize}
  \item We formalize three strengths of typed patchwork
    (\S\ref{sec:prelim}, \S\ref{sec:typed}) and separate them from
    the untyped Renz--Nebel property.
  \item We attempt three candidate counterexamples
    (\S\ref{sec:probe1}--\S\ref{sec:probe3}) selected for structural
    diversity: forced coincidence of distinct types under shared
    parentage, inverse-role cycle closure, and deep witness chains
    with interleaved universals.
  \item We analyse each probe and report that each fails to produce a
    counterexample, with the obstruction reducing either to local
    arc-inconsistency or to type-level unsatisfiability
    (\S\ref{sec:analysis}).
  \item We identify the precise structural shape a genuine typed-patchwork
    failure would need to take (\S\ref{sec:shape}), and argue it
    aligns with the dual of the completeness-extraction key lemma.
  \item We conclude with an honest assessment of what this probe does
    and does not establish (\S\ref{sec:honest}).
\end{itemize}

This paper is deliberately short and deliberately negative.  It is a
research probe, not a proof.  We invite the reader to find the
counterexample we could not.

\section{Preliminaries}
\label{sec:prelim}

We briefly recap $\mathsf{RCC5}$, Hintikka types, and the Renz--Nebel
patchwork property.  The reader is referred
to~\cite{OverviewClaude, SplitForest, CompletenessExtraction} for
full definitions.

\paragraph{$\mathsf{RCC5}$.}  The five base relations are
$\DR, \PO, \EQ, \PP, \PPI$, with $\EQ$ interpreted as identity
(strong-$\EQ$ semantics) and $\PP, \PPI$ as converse partial orders.
For distinct nodes, the four admissible atomic relations are
$\RCC$.  The composition table $\comp$ is well known; we write
$\comp(R_1, R_2)$ for the set of $\mathsf{RCC5}$ relations consistent
with the composition of $R_1$ and $R_2$.

\paragraph{Atomic networks and patchwork.}
An \emph{atomic $\mathsf{RCC5}$ network} is a pair $(V, \rho)$ where
$V$ is finite and $\rho: V \times V \to \RCCall$ satisfies
$\rho(x,x) = \EQ$, $\rho(y,x) = \inv(\rho(x,y))$, and $\rho$ is
\emph{path-consistent}:
\[
  \rho(x,z) \in \comp(\rho(x,y), \rho(y,z))
  \quad \text{for all } x, y, z \in V.
\]

\begin{theorem}[Renz and Nebel~\cite{RenzNebel1999}]
\label{thm:rn}
A finite atomic $\mathsf{RCC5}$ network is realizable iff it is
path-consistent.
\end{theorem}

\paragraph{Hintikka types.}
Fix a concept $C_0$ and a TBox $T$; let $\cl(C_0, T)$ be the
closure under subconcepts, negations, and TBox axioms.  A
\emph{Hintikka type} (for $(C_0, T)$) is a subset
$\tau \subseteq \cl(C_0, T)$ that is maximally locally consistent:
for every Boolean subconcept it contains exactly one of
$\{D, \neg D\}$, for each $D_1 \sqcap D_2$ it contains $D_1$ and
$D_2$, for each $D_1 \sqcup D_2$ it contains at least one, and so on.
We write $\Tp(C_0, T)$ for the finite set of Hintikka types.

\paragraph{Type-safety.}
Given an edge $\rho(x, y) = R$ with $\tp(x) = \alpha$, $\tp(y) = \beta$:
the edge is \emph{type-safe} if for every universal $\forall S.D
\in \alpha$ with $S = R$, we have $D \in \beta$, and symmetrically
for $\inv(R)$.  Write $\TypeSafe(\alpha, \beta) \subseteq \RCC$ for
the set of $\mathsf{RCC5}$ relations $R$ with $(\alpha, R, \beta)$
type-safe.

\section{Three strengths of typed patchwork}
\label{sec:typed}

We now state three progressively stronger patchwork claims.  The
first is Renz--Nebel; the third is what the decidability proof needs.

\begin{definition}[Atomic patchwork]
\label{def:atomic-pw}
For every finite path-consistent atomic $\mathsf{RCC5}$ network
$(V, \rho)$, there exists a model $M$ of the $\mathsf{RCC5}$ algebra
and an injection $\iota: V \to M$ such that
$\rho(x,y) = \rho^M(\iota(x), \iota(y))$ for all $x,y \in V$.
\end{definition}

\begin{theorem}
Atomic patchwork holds (Theorem~\ref{thm:rn}).
\end{theorem}

\begin{definition}[Locally-coherent typed patchwork]
\label{def:lc-pw}
Fix an $\ALCIRCC{5}$ TBox $T$ and concept $C_0$.  A \emph{typed
network} is a triple $(V, \rho, \tau)$ with $(V, \rho)$ a
path-consistent atomic $\mathsf{RCC5}$ network and
$\tau : V \to \Tp(C_0, T)$ such that every edge is type-safe.
Locally-coherent typed patchwork says: every typed network extends
to a model $M \models T \cup \{C_0\}$ with an injection
$\iota: V \to M$ preserving $\rho$ and $\tau$.
\end{definition}

This is already stronger than Definition~\ref{def:atomic-pw}: it
requires not merely an $\mathsf{RCC5}$ model but a model that
satisfies the TBox with the prescribed types.  Additional existential
witnesses must be supplied for every $\exists R.D \in \tau(x)$, and
those witnesses must be globally consistent with the rest of $V$.

\begin{definition}[Arc-consistent typed patchwork]
\label{def:ac-pw}
Given $T$, $C_0$ as above.  A \emph{disjunctive typed network} is
$(V, D, \tau)$ where $D: V \times V \to 2^\RCC$ assigns a non-empty
set of admissible atomic relations to each ordered pair, subject
to $\inv$-symmetry.  \emph{Arc-consistency} is the fixed point of
the propagation
\[
  D'(x, z) \;:=\; D(x, z) \,\cap\,
  \bigcup_{R_1 \in D(x,y)} \comp(R_1, D(y, z))
  \,\cap\, \TypeSafe(\tau(x), \tau(z)).
\]
Arc-consistent typed patchwork says: if $(V, D, \tau)$ is
arc-consistent and every $D(x,y)$ is non-empty, then there exists
a model $M \models T \cup \{C_0\}$ realizing some atomic
refinement of $(V, D, \tau)$.
\end{definition}

The split-forest semantics paper~\cite{SplitForest} establishes
arc-consistent typed patchwork for a \emph{restricted} class of
disjunctive networks (those arising from valid finite rank-$d$
quotients) via the LCA-induction proof of Theorem~1.17 plus
König's lemma on canonical refinements.  The completeness-extraction
paper~\cite{CompletenessExtraction} establishes the converse: every
model induces a quotient that satisfies arc-consistency, via the key
lemma that realized relations are automatically relation-safe.

\paragraph{The gap.}  What is \emph{not} established, to our
knowledge, is Definition~\ref{def:ac-pw} for \emph{arbitrary}
arc-consistent typed networks---not just those arising from quotients
of genuine models.  This is where a counterexample could live.  If
there is an arc-consistent $(V, D, \tau)$ not realizable by any
model of $T$, then arc-consistency alone is not complete, and the
decidability argument needs further propagation.

We turn to three probes for such a counterexample.

\section{Probe 1: forced coincidence of distinct types}
\label{sec:probe1}

The idea: force three nodes $x, y, z$ with $x \PP y$ and
$x \PP z$, and use the TBox to demand that $y$ and $z$ have types
$B$ and $C$ that are pairwise incompatible on every non-$\EQ$
atomic relation.  Renz--Nebel allows $y ? z \in \{\DR, \PO, \EQ,
\PP, \PPI\}$; if the TBox forbids four of the five via
$\TypeSafe$, only $\EQ$ remains, and if $B \sqcap C$ is itself
unsatisfiable as a type, the network is unrealizable.

\begin{probe}[Forced-coincidence probe]
\label{probe:fc}
Let $A, B, C$ be distinct concept names.  Let
\begin{align*}
  A &\sqsubseteq \exists \PP. B \sqcap \exists \PP. C, \\
  B &\sqsubseteq \forall \DR. \neg C
      \sqcap \forall \PO. \neg C
      \sqcap \forall \PP. \neg C
      \sqcap \forall \PPI. \neg C, \\
  B &\sqsubseteq \neg C.
\end{align*}
Take $V = \{x, y, z\}$ with $\tau(x) \supseteq \{A\}$,
$\tau(y) \supseteq \{B\}$, $\tau(z) \supseteq \{C\}$, and
$\rho(x,y) = \PP$, $\rho(x,z) = \PP$.
\end{probe}

\begin{observation}
\label{obs:fc-fail}
The network in Probe~\ref{probe:fc} is \emph{not} arc-consistent.
\end{observation}

\begin{proof}
The propagation $D(y,z)$ starts from
$\comp(\PPI, \PP) = \RCCall$ (path-consistency allows all five).
Type-safety filters by $\tau(y) = B$: the universals in $B$ rule out
all non-$\EQ$ relations on the $y$-side, so
$\TypeSafe(\tp(y), \tp(z)) \cap \comp(\PPI, \PP) \subseteq \{\EQ\}$.
But $\rho(y, z) = \EQ$ forces $y = z$, hence $\tp(y) = \tp(z)$,
hence $B \sqcap C \in \tp(y)$---but $B \sqsubseteq \neg C$, so
$B \sqcap C$ is type-inconsistent and there is no Hintikka type
containing it.  Arc-consistency propagation therefore empties
$D(y, z)$.
\end{proof}

The probe reduces to an \emph{arc-inconsistency}.  It is caught by
the AC-3 loop and never presented to the global realizability
question at all.  This is exactly what one would hope: typed
patchwork is not even tested here, because the local filter
already fails.

The probe is structurally isomorphic to the ``coincidence
obstruction'' of Wessel 2002~\cite[Fig.~11]{Report7} \emph{except
that in Wessel's setting, $\EQ$ is not ruled out by a universal};
there, the obstruction is that nothing forces the cells to be
distinct, so the grid cells silently identify.  In our probe, the
TBox \emph{forces} the cells to be distinct via $B \sqsubseteq
\neg C$, but then arc-consistency catches the conflict locally.
The two obstructions are dual.

\section{Probe 2: inverse-role cycle closure}
\label{sec:probe2}

The idea: use inverse roles to force a cycle
$x \to y \to z \to x$ of role compositions, and hope that the
composition table permits the cycle locally at each arc but no
global atomic assignment closes it consistently with the TBox.

\begin{probe}[Inverse-cycle probe]
\label{probe:ic}
Let $A, B, C, D$ be concept names.  Consider
\begin{align*}
  A &\sqsubseteq \exists \PP. B, \\
  B &\sqsubseteq \exists \PO. C, \\
  C &\sqsubseteq \exists \PPI. D, \\
  D &\sqsubseteq \forall \PP. \neg A.
\end{align*}
Take $V = \{x, y, z, w\}$ with
$\tau(x) \supseteq \{A\}$, $\tau(y) \supseteq \{B\}$,
$\tau(z) \supseteq \{C\}$, $\tau(w) \supseteq \{D\}$, and edges
$\rho(x,y) = \PP$, $\rho(y,z) = \PO$, $\rho(z,w) = \PPI$.
The fourth edge $\rho(x, w)$ is unconstrained by explicit
existentials.
\end{probe}

\begin{observation}
\label{obs:ic-ok}
The network in Probe~\ref{probe:ic} is arc-consistent and
realizable.
\end{observation}

\begin{proof}
Compose the path $x \to y \to z \to w$:
$\comp(\PP, \PO) \ni \PP, \PO, \DR$; then with $\PPI$ on the
$z$-$w$ edge, we get $\rho(x, w) \in \comp(\comp(\PP, \PO), \PPI)$,
which is all of $\RCCall$ (the composition of
$\{\PP, \PO, \DR\}$ with $\PPI$ covers every relation).
Type-safety from $D \sqsubseteq \forall \PP. \neg A$ rules out
$\rho(w, x) = \PP$, i.e., $\rho(x, w) = \PPI$.  Remaining
candidates: $\{\DR, \PO, \PP, \EQ\}$.  $\EQ$ is ruled out if
$A \sqcap D$ is type-inconsistent (it is, via the universal on $D$
reflecting back: if $w = x$ then $w \PP \cdot$ but $x \PP y$ with
$\tau(y) = B \neq \neg A$---wait, this doesn't immediately fire
unless $A$ is in $\tau(y)$, which it is not).  Actually $\EQ$
merely forces $\tau(x) = \tau(w)$, hence $\{A, D\} \subseteq
\tau$.  This is type-consistent only if $A \sqcap D$ is
satisfiable, which is not ruled out by the TBox.  We pick
$\rho(x, w) = \DR$ for a concrete realization: a model with four
distinct points and the specified relations satisfies all axioms.
\end{proof}

The probe produces an arc-consistent, realizable network.  No
obstruction.

\paragraph{Variants.}  We attempted several variants: longer
cycles (5, 6 nodes), cycles with $\EQ$-forcing universals,
cycles that interleave $\PO$ and $\PP$ in ways designed to
exhaust the composition table's slack.  Every variant either
(a) produced a realizable network, or (b) collapsed to arc-inconsistency
via type-safety filtering on some edge.  We did not find a variant
that produced a path-consistent, arc-consistent, type-coherent
network without a model.

\section{Probe 3: deep witness chains with interleaved universals}
\label{sec:probe3}

The idea: an $\mathcal{ALC}$-style ``yardstick'' construction.  Use
a long chain of existential witnesses with universals that accumulate
along the chain, forcing an obligation that cannot be satisfied at
the chain's end without backtracking.

\begin{probe}[Yardstick probe]
\label{probe:ys}
Fix $n \geq 2$ and concept names $A_0, A_1, \ldots, A_n, F$.  Let
\begin{align*}
  A_i &\sqsubseteq \exists \PP. A_{i+1} \quad (0 \leq i < n), \\
  A_i &\sqsubseteq \forall \PP. F \quad (0 \leq i \leq n-1), \\
  A_n &\sqsubseteq \neg F.
\end{align*}
\end{probe}

\begin{observation}
\label{obs:ys-type}
Probe~\ref{probe:ys} is type-unsatisfiable for every $n \geq 1$.
\end{observation}

\begin{proof}
Every model of $A_0$ has a $\PP$-successor of type $A_1$; by the
universal $\forall \PP. F \in A_0$, this successor also has $F$.
But $A_1$ contains $\forall \PP. F$ and (for $n \geq 2$) by
iteration $A_{n-1}$ produces $A_n$ with $F$, contradicting
$A_n \sqsubseteq \neg F$.  So there is no Hintikka type containing
$A_0$.
\end{proof}

\paragraph{Refining the probe.}  The above probe is caught by
type-saturation before any network is constructed.  A more
interesting variant replaces the universal with something that
only fires conditionally on another concept that might be absent
at the chain's end.  We tried:
\begin{align*}
  A_i &\sqsubseteq \exists \PP. A_{i+1}, \\
  A_i &\sqsubseteq \forall \PP. (G \sqcup \neg G), \\
  A_n &\sqsubseteq \neg G \sqcap \neg (\neg G).
\end{align*}
Of course $\forall \PP.(G \sqcup \neg G)$ is trivially true and the
bottom constraint $A_n \sqsubseteq \neg G \sqcap G$ is vacuously
unsatisfiable in types.  After several variants, the failure
mode was uniform: either the type closure caught an inconsistency
immediately, or the type was satisfiable and a trivial model
existed.

The underlying reason: $\mathcal{ALC}$ universals propagate
deterministically along role chains.  In $\mathcal{ALCI}$, inverse
roles are also deterministic.  In $\ALCIRCC{5}$ on top, the
role labels are further constrained by $\mathsf{RCC5}$ composition
via arc-consistency, but composition only \emph{loosens} the
constraint (unions of allowed relations), never tightens it below
what the types already force.  The yardstick construction is
therefore flattened by arc-consistency rather than amplified.

\section{Structural analysis: why every probe fails}
\label{sec:analysis}

Across all three probes (and several unreported variants), the
obstruction to realizability reduced to one of two phenomena:
\begin{enumerate}
  \item \textbf{Local arc-inconsistency:}
    $\TypeSafe(\tau(x), \tau(y)) \cap \comp(\ldots) = \varnothing$
    for some pair, and AC-3 empties $D(x,y)$.
  \item \textbf{Type-level unsatisfiability:}
    $\bigcap_{\forall R.D_i \text{ firing at } \tau(x)} \{D_i\}$
    is inconsistent, and the Hintikka-type closure catches it.
\end{enumerate}
Neither of these is a failure of typed patchwork in the sense of
Definition~\ref{def:ac-pw}, because both occur \emph{before}
arc-consistency terminates with all $D(x,y)$ non-empty.

\paragraph{A structural reason.}
The completeness-extraction paper~\cite{CompletenessExtraction}
identifies the key lemma: in a genuine model $M$, every realized
atomic relation between two elements $x, y \in M$ is automatically
relation-safe for $(\tp^M(x), \tp^M(y))$.  The dual statement is:
if no atomic relation is relation-safe for $(\alpha, \beta)$
under composition $\comp(R_1, R_2)$, then no model can realize
$\alpha$ and $\beta$ at those positions.

Typed patchwork is essentially the contrapositive: \emph{if} every
arc of a typed network has a non-empty filtered domain, \emph{then}
some model realizes the network.  The probes fail to produce a
counterexample because, intuitively, non-empty filtered domains
\emph{are} the witness that a model exists.

\begin{proposition}[Conjectural; we were unable to falsify]
\label{prop:conjecture}
For every $\ALCIRCC{5}$ TBox $T$ and satisfiable concept $C_0$,
typed patchwork holds in the sense of
Definition~\ref{def:ac-pw}: every arc-consistent typed network
with non-empty filtered domains is realizable by some model of
$T \cup \{C_0\}$.
\end{proposition}

A full proof of Proposition~\ref{prop:conjecture} would amount to
a simpler alternative to the split-forest construction: extend the
network by Henkin witnesses, maintaining arc-consistency at each
step.  The step that is \emph{not obviously correct} is that
adding a new witness for some existential $\exists R.D$ at a node
$x$ requires choosing $\rho(y, z)$ for every existing $z \in V$,
and these choices must be simultaneously arc-consistent.  The
split-forest paper handles this via canonical refinements and
König's lemma on an infinite tree of choices.  We do not prove
Proposition~\ref{prop:conjecture} here; we merely note that our
probes were consistent with it.

\section{The shape of a genuine failure}
\label{sec:shape}

Suppose, contrary to Proposition~\ref{prop:conjecture}, that
typed patchwork fails.  What would the counterexample look like?

\begin{enumerate}
  \item It would be an arc-consistent typed network with all
    $D(x, y) \neq \varnothing$.
  \item The model-realizability question would fail \emph{not}
    at any individual arc (all arcs are locally consistent), but
    at a \emph{global} level where the simultaneous atomic
    choices cannot be made.
  \item The failure must survive arc-consistency.  Since AC-3
    propagates through ternary composition, the failure
    mechanism must be essentially $k$-ary for some $k \geq 4$:
    four or more nodes whose pairwise filtered domains are all
    non-empty but whose joint atomic assignments are all
    inconsistent with the TBox.
\end{enumerate}

\paragraph{Why this shape is hard to achieve.}  Renz--Nebel
Theorem~\ref{thm:rn} already forecloses the $\mathsf{RCC5}$-only
version of such $k$-ary failures: atomic path-consistency is
complete for realizability of untyped networks.  To manufacture a
$k$-ary typed failure, one would need the TBox to introduce a
$k$-ary constraint that is \emph{invisible to any triple of nodes}
but \emph{visible to the full $k$-subset}.

In $\mathcal{ALCI}$, every constraint is expressible as a
restriction on a single node's type (universals, existentials,
Boolean combinations).  The TBox does not directly express
$k$-ary constraints for $k \geq 2$.  TBox axioms become
node-local (unary) constraints after Hintikka closure.  Arc
constraints ($k = 2$) are captured by type-safety.  Triple
constraints ($k = 3$) are captured by path-consistency /
composition.  \emph{There is no $\mathcal{ALCI}$ mechanism for
introducing a genuinely $k$-ary constraint for $k \geq 4$.}

This is the structural reason we conjecture typed patchwork holds.
Any $k$-ary-looking constraint in $\mathcal{ALCI}$ decomposes into
unary, binary (type-safety), or ternary (composition) pieces, all
of which are captured by arc-consistency.

\paragraph{What would break this reasoning.}  The only loophole we
see is if \emph{inverse roles plus composition} could
\emph{synthesize} a $k$-ary constraint through a long chain.  The
canonical example is $\exists R. \forall R^{-}. A$: this uses both
directions of a single role to express ``every $R$-predecessor of
some $R$-successor is $A$.''  In principle, iterating such
constructions could accumulate global constraints that no finite
arc-consistency check sees.

However, in the $\ALCIRCC{5}$ context, role inverse is not
independent of $\mathsf{RCC5}$ composition: $\mathsf{RCC5}$
relations are their own inverse under $\inv$, and composition is
associative.  Long chains of $\exists R. \forall R^-. \ldots$
collapse to unary constraints under type closure.  So we do not
see a route to $k$-ary synthesis via inverse role alone.

If such a route exists, the counterexample to Proposition~\ref{prop:conjecture}
would be the vehicle to exhibit it, and it would simultaneously be
a counterexample to the decidability proof.  This is the single
tightest place to look for undecidability.

\section{Implications for undecidability}
\label{sec:undec}

If typed patchwork holds (Proposition~\ref{prop:conjecture}):
\begin{itemize}
  \item The decidability proof for $\ALCIRCC{5}$ is robust.
  \item Every standard undecidability reduction (grid encodings,
    PCP, word problems) is blocked by the same mechanism: such
    reductions require global rigidity, which typed arc-consistency
    does not permit.
  \item A non-trivial undecidability proof would need to bypass
    patchwork entirely---encode computation into something other
    than the relational structure.  Given that
    $\ALCIRCC{5}$'s only semantic content \emph{is} the
    relational structure, such a bypass seems unavailable.
\end{itemize}

If typed patchwork fails:
\begin{itemize}
  \item The failure exhibits a concrete $k$-ary constraint
    ($k \geq 4$) that arc-consistency misses.
  \item One would then need to amplify the failure into a
    computation-simulating reduction.  This is non-trivial but not
    impossible: a single $k$-ary failure at fixed $k$ does not
    imply undecidability, but a \emph{family} of failures indexed
    by the tape content of a Turing machine would.
\end{itemize}

We emphasize: \emph{finding a typed-patchwork failure does not by
itself prove undecidability.}  It only opens the door.  But
\emph{not} finding one---as in this paper---closes the most obvious
door we had.

\section{Honest assessment}
\label{sec:honest}

This paper is a negative result in both directions.  We did not
find a counterexample, and we did not prove
Proposition~\ref{prop:conjecture}.  What we did establish:

\begin{enumerate}
  \item The typed patchwork question can be precisely formalised
    (Definition~\ref{def:ac-pw}) as a question about arc-consistency
    completeness.
  \item Several natural candidate counterexamples
    (forced coincidence, inverse cycle, yardstick) reduce to
    arc-inconsistency or type-inconsistency before reaching the
    global realizability question.
  \item The structural reason for this reduction is that
    $\mathcal{ALCI}$ cannot express $k$-ary constraints for
    $k \geq 4$ that are invisible to arc-consistency.
\end{enumerate}

What we did \emph{not} establish:

\begin{enumerate}
  \item A full proof of typed patchwork.  This would require
    Henkin-style model construction maintaining arc-consistency,
    essentially a repackaging of the split-forest soundness
    direction.
  \item A sharp characterisation of which $\mathcal{ALCI}$
    constructions could in principle synthesise $k$-ary constraints
    through inverse-role interaction.  We sketched an argument that
    none do, but did not prove it rigorously.
  \item An exhaustive search of candidate counterexamples.  Three
    probes and a handful of variants is not exhaustive; a reader
    with a better combinatorial imagination may well find one.
\end{enumerate}

\paragraph{Outlook.}  The most promising refinement of this probe
would be to formalize the ``no $k$-ary synthesis'' claim and prove
it rigorously.  A theorem of the form

\begin{quote}\emph{
  Every $\mathcal{ALCI}_{\mathsf{RCC5}}$ constraint on a finite
  network decomposes into constraints of arity at most three, and
  arc-consistency is complete for all arity-$\leq 3$ constraint
  satisfaction problems over $\mathsf{RCC5}$.
}\end{quote}

\noindent would close both questions simultaneously: it would
prove typed patchwork and simultaneously certify that no
$\ALCIRCC{5}$ TBox can express an undecidable CSP.

A theorem of that shape would provide an alternative, more direct
decidability proof than the split-forest construction, grounded in
constraint satisfaction rather than tree semantics.  It would also
finally give a \emph{principled} answer to why every undecidability
reduction fails: not a survey of blocked techniques, but a single
theorem that all are instances of.

We leave this to future work.

\begin{thebibliography}{99}

\bibitem{RenzNebel1999}
J.~Renz and B.~Nebel.
\newblock On the complexity of qualitative spatial reasoning: A maximal
  tractable fragment of the region connection calculus.
\newblock \emph{Artificial Intelligence}, 108(1--2):69--123, 1999.

\bibitem{SplitForest}
GPT-5.4 (OpenAI), prompted by M.~Wessel.
\newblock Sibling-interface descriptors and split-forest semantics for
  $\ALCIRCC{5}$.
\newblock Manuscript,
  \texttt{papers/trees/sibling\_interface\_descriptors\_ALCIRCC5\_completed\_eqsync\_canonical\_needpatched.pdf},
  April 2026.

\bibitem{CompletenessExtraction}
Claude (Anthropic, Opus 4.6), prompted by M.~Wessel.
\newblock Completeness extraction for $\ALCIRCC{5}$.
\newblock Manuscript,
  \texttt{papers/completeness\_extraction\_ALCIRCC5.pdf}, April 2026.

\bibitem{OverviewClaude}
Claude (Anthropic, Opus 4.7), prompted by M.~Wessel.
\newblock Overview of the $\ALCIRCC{5}$ decidability project.
\newblock Manuscript, \texttt{papers/overview\_ALCIRCC5.pdf}, April 2026.

\bibitem{Report7}
M.~Wessel.
\newblock Report on qualitative spatial reasoning with description logics.
\newblock Technical Report, 2002--2003.
\newblock \texttt{papers/report7.pdf}.

\bibitem{README}
M.~Wessel et al.
\newblock Repository README: the $\ALCIRCC{5}$ decidability project.
\newblock \texttt{github.com/lambdamikel/alcircc5}, 2026.

\end{thebibliography}

\end{document}
